%% kapa: exam=1 round=20 year=2009 period=1 form=A subject=economics subject_ko=경제원론 seq=34 number=74
%% source: Q-Net 감정평가사 기출문제 (cst00309) artlSeq=5024178 file=extracted/1008/5024178/2009년 제20회 감정평가사 1차시험 기출문제/1교시.hwp
%% conversion: kordoc 4.24.51 -> tools/gen1.py (auto); review=proofread-sonnet; flags=none
\begin{question}{74}
A는 사과주스 한 잔과 당근주스 한 잔을 서로 바꾸어 마셔도 동일한 만족을 얻는 반면, B는 사과주스와 당근주스는 반드시 1:1로 섞어 마셔야 만족한다. 다음 설명 중 옳은 것은?
\begin{choices}
  \item A에게 당근주스와 사과주스는 완전보완재이다.
  \item A에게 당근주스와 사과주스의 한계대체율은 1이다.
  \item A의 무차별곡선은 수평선이다.
  \item B의 무차별곡선은 수평선이다.
  \item B에게 당근주스와 사과주스의 한계대체율은 -1이다.
\end{choices}
\end{question}
\answer{2}
